% !PRINTED 279
\weberpage{279}
Nun l\"asst sich andererseits wieder $\omega$ so klein annehmen, dass die Wurzeln von $L_n' = 0$ von denen von $L_n = 0$ beliebig wenig unterschieden sind (\S.\,40).

Es kann zwar eine Doppelwurzel von $L_n=0$ in zwei einfache Wurzeln von $L_n'=0$ \"ubergehen; aber da $L_n'=0$ keine imagin\"aren Wurzeln hat, so sind diese reell; und dasselbe findet statt, wenn $L_n=0$ mehrfache Wurzeln hat.

\gesperrt{Hiernach behalten die S\"atze I, II ihre G\"ultigkeit, wenn die Voraussetzung aufgegeben wird, dass in der Kette der $L_k$ keine auf einander folgenden Glieder verschwinden; nur m\"ussen die mehrfachen Wurzeln dabei nach ihrer Vielfachheit gez\"ahlt werden.}

\begin{center}
\S.\,89.

\gesperrt{Die Sturm'schen Functionen.}
\end{center}

Nach diesen besonderen Beispielen wenden wir uns zur Betrachtung des Verfahrens, durch das man in allen F\"allen eine Sturm'sche Kette erh\"alt\footnote{\gesperrt{Sturm}, M\'em. sur la r\'esolution des \'equations num\'eriques. M\'em. de l'acad\'emie de Paris. Sav. \'etrang. VI, 1835. Auszug in Bull. de Ferussac XI, 1829.}. Die Bildungsweise dieser Functionen ist principiell ausserordentlich einfach, wenn auch in der praktischen Ausf\"uhrung meist nicht durchf\"uhrbar.

Wir beschr\"anken uns hier auf die Betrachtung von Gleichungen \gesperrt{ohne mehrfache Wurzeln}, oder wir nehmen an, dass vor der Anwendung $f(x)$ von jedem gemeinschaftlichen Factor mit seiner Derivirten $f'(x)$ befreit sei.

Wenn wir dann
\weq{1}{f_1(x) = f'(x)}
annehmen, so ist sicher die Bedingung \S.\,86, 4. befriedigt. Nun verfahren wir so, als ob es sich um die Aufsuchung des gr\"ossten gemeinschaftlichen Theilers von $f(x)$ und $f_1(x)$ handle, indem wir dabei jedesmal das Vorzeichen des Restes umkehren; wir bilden also durch Division die Gleichungen
\weq{2}{\begin{gathered}
f = q_1 f_1 - f_2\\
f_1 = q_2 f_2 - f_3\\
\cdots\cdots\\
f_{m-2} = q_{m-1} f_{m-1} - f_m,
\end{gathered}}
%CONT>>
